% Part: computability
% Chapter: computability-theory
% Section: rice-theorem

\documentclass[../../../include/open-logic-section]{subfiles}

\begin{document}

\olfileid{cmp}{thy}{rce}
\olsection{Teorema Rice}

Yen dipikirake, kowe bakal weruh yen sipat-sipat mligi
$\fn{Tot}$ ora melu nemtokake bukti \olref[tot]{prop:total}.
Kita ngrancang $h(x,y)$ supaya tumindake padha karo
fungsi konstan $j(y) = 0$ persis nalika $x$ kalebu $K$;
nanging ing kahanan iku kita uga bisa ngrancang supaya
tumindake padha karo fungsi komputabel parsial liyane
sembarang. Pangamatan iki ngidini kita nyebut teorema
kang luwih umum: kurang luwih, ora ana sipat ora trivial
saka fungsi komputabel kang bisa diputusake.

Elinga yen $\cfind{0}$, $\cfind{1}$, $\cfind{2}$,~\dots
minangka enumerasi baku fungsi-fungsi komputabel parsial.

\begin{thm}[Teorema Rice]
  Tetepna $C$ minangka himpunan sembarang fungsi
  komputabel parsial, lan tetepna $A =
  \Setabs{n}{\cfind{n} \in C}$. Yen $A$ komputabel,
  mula $C$ kosong utawa $C$ minangka himpunan kabeh
  fungsi komputabel parsial.
\end{thm}

{\em Himpunan indeks} yaiku himpunan $A$ kang nduwé sipat
mangkene: yen $n$ lan $m$ indeks kang ``ngitung'' fungsi
kang padha, mula $n$ lan $m$ loro-lorone kalebu $A$,
utawa loro-lorone ora kalebu. Gampang dideleng yen
himpunan~$A$ ing teorema iki nduwé sipat mau. Kosok
baline, yen $A$~himpunan indeks lan $C$~himpunan fungsi
kang diitung dening indeks-indeks iku, mula
$A = \Setabs{n}{\cfind{n} \in C}$.

\begin{explain}
Kanthi istilah iki, teorema Rice padha tegese karo
pratelan yen ora ana himpunan indeks ora trivial kang
bisa diputusake. Kanggo mangerteni isi teorema iki,
migunani yen kita mbedakake \emph{program}
(umpamane, ing basa pamrograman kang koksenengi)
saka fungsi kang diitung dening program iku.
Mesthi ana pitakon bab program (indeks), yaiku
obyek sintaktis, kang bisa diwangsuli sacara komputabel:
apa program iki nduwé luwih saka 150 simbol?
Apa dawane luwih saka 22 larik? Apa ana pratelan
``while'' ing jerone? Apa rerangken aksara ``hello world''
tau dadi argumen kanggo pratelan ``print'' ing teks program?
Teorema Rice nyebut yen ora ana pitakon ora trivial
bab \emph{tumindake} program kang bisa diwangsuli
sacara komputabel. Cathetan panyunting: tumindak ing teorema iki sipat fungsi parsial kang diitung, kalebu domain lan asil, dudu wektu mlaku utawa jejak komputasi kang ora ditemtokake dening fungsi kasebut. Iki kalebu pitakon kaya mangkene:
apa program mandheg nalika nampa input $0$?
Apa program tau mandheg? Apa program tau
ngasilake wilangan genap?
\end{explain}

\begin{proof}[Bukti teorema Rice]
Upamane $C$ ora kosong lan uga dudu himpunan kabeh
fungsi komputabel parsial, lan tetepna $A$ minangka
himpunan indeks fungsi ing~$C$. Kita bakal nuduhake
yen $A$ komputabel, kita bisa ngrampungake masalah
mandheg; mula $A$~ora komputabel.

Tanpa ngurangi umumé argumen, kita bisa nganggep
yen fungsi $f$ kang ora ditegesi ing input endi wae
ora kalebu $C$ (yen kosok baline, ijolna $C$ lan
komplemene ing argumen ing ngisor iki). Tetepna $g$
minangka fungsi sembarang ing~$C$. Gagasane yaiku
yen kita bisa mutusake~$A$, kita bisa mbedakake
indeks kang ngitung~$f$ saka indeks kang ngitung~$g$;
banjur kemampuan iku bisa digunakake kanggo
ngrampungake masalah mandheg.

Carane mangkene. Kanthi nggunakake fungsi-fungsi
komputabel parsial universal, kita bisa netepake fungsi
\[
h(x,y) \simeq
\begin{cases}
\text{ora ditegesi} & \text{yen $\cfind{x}(x) \fundefined$} \\
g(y) & \text{yen ora.}
\end{cases}
\]
Kanggo ngitung $h$, luwih dhisik kita nyoba ngitung
$\cfind{x}(x)$; yen komputasi iku mandheg, kita
nerusake ngitung~$g(y)$; lan yen komputasi
{\em iku} uga mandheg, kita mbalekake asile.
Luwih formal, kita bisa nulis
\[
h(x,y) \simeq \Proj{2}{0}(g(y),\fn{Un}(x,x)).
\]
ing ngendi $\Proj{2}{0}(z_0, z_1) = z_0$ minangka
fungsi proyeksi $2$ argumen kang mbalekake argumen
kapangka-$0$, lan fungsi iki komputabel.

Mula $h$ minangka komposisi fungsi-fungsi komputabel
parsial, lan sisih tengen ditegesi lan padha
karo~$g(y)$ yen lan mung yen $\fn{Un}(x,x)$ lan $g(y)$
loro-lorone ditegesi.

Gatekna yen kanggo~$x$ kang tetep, yen
$\cfind{x}(x)$ ora ditegesi, mula $h(x,y)$
ora ditegesi kanggo saben~$y$; lan yen
$\cfind{x}(x)$ ditegesi, mula $h(x,y) \simeq g(y)$.
Dadi, kanggo saben nilai~$x$ kang tetep,
$h(x,y)$ tumindake padha karo $f$ utawa padha
karo $g$, lan mutusake apa $\cfind{x}(x)$
ditegesi padha karo mutusake endi saka loro
kahanan iku kang kelakon. Iki uga padha karo
mutusake apa fungsi $h_x(y) \simeq h(x,y)$ kalebu~$C$,
lan yen $A$ komputabel, kita bisa nindakake iku.

Luwih formal, amarga $h$~komputabel parsial,
fungsi iku padha karo $\cfind{e}$ kanggo sawijining
indeks~$e$. Cathetan koreksi sumber OLPL-768: amarga fungsi iki nduweni rong argumen, indeks kasebut kudu diwaca kanggo varian aritas loro utawa liwat kode pasangan. Miturut teorema $s$-$m$-$n$, ana
fungsi rekursif primitif~$s$ saengga kanggo saben~$x$,
$\cfind{s(e,x)}(y) = h_x(y)$. Saiki kita nduwé yen
kanggo saben $x$, yen $\cfind{x}(x) \fdefined$,
mula $\cfind{s(e,x)}$ fungsi kang padha karo~$g$,
lan mula $s(e,x)$ kalebu $A$. Kosok baline,
yen $\cfind{x}(x) \uparrow$, mula
$\cfind{s(e,x)}$ fungsi kang padha karo~$f$,
lan mula $s(e,x)$ ora kalebu~$A$. Kanthi tembung
liya, kanggo saben~$x$, $x
\in K$ yen lan mung yen $s(e,x) \in A$.
Yen $A$ komputabel, $K$~uga komputabel,
kang nuwuhake kontradiksi. Dadi $A$ ora komputabel.
\end{proof}

Teorema Rice migunani banget. Korolari langsung
ing ngisor iki nuduhake sawetara tuladha panganggone.
\begin{cor}
Himpunan-himpunan ing ngisor iki ora bisa diputusake.
\begin{enumerate}
\item $\Setabs{x}{\text{$17$ kalebu citra $\cfind{x}$}}$
\item $\Setabs{x}{\text{$\cfind{x}$ konstan}}$
\item $\Setabs{x}{\text{$\cfind{x}$ total}}$
\item $\Setabs{x}{\text{saben $y < y'$, $\cfind{x}(y) \fdefined$, lan
    yen $\cfind{x}(y') \fdefined$, mula $\cfind{x}(y) < \cfind{x}(y')$}}$
\end{enumerate}
\end{cor}

\begin{proof}
  Kabeh iku himpunan indeks kang ora trivial.
\end{proof}

\end{document}
